% ======================================================================
%  Kapustin–Witten 理论导读讲义 — 共享导言区
%  编译：xelatex main.tex（两遍），工作目录 lecture_notes/
% ======================================================================

% ---------- 数学 ----------
\usepackage{amsmath,amssymb,amsthm,mathtools}
\usepackage{mathrsfs}
\usepackage{bm}

% ---------- 排版 ----------
\usepackage[margin=2.35cm,top=2.5cm,bottom=2.5cm]{geometry}
\usepackage{graphicx}
\usepackage{booktabs,array,longtable,multicol,multirow}
\usepackage{enumitem}
\setlist{itemsep=2pt,topsep=3pt,parsep=0pt}
\usepackage{xcolor}
\usepackage{xspace}
\usepackage{fancyhdr}
\usepackage[colorlinks=true,linkcolor=blue!55!black,urlcolor=teal!70!black,citecolor=blue!55!black]{hyperref}

% ---------- 图 ----------
\usepackage{tikz}
\usetikzlibrary{arrows.meta,positioning,calc,shapes.geometric,fit,backgrounds,decorations.pathreplacing}
\usepackage{tikz-cd}

% ---------- 颜色 ----------
\definecolor{kwblue}{RGB}{28,66,120}
\definecolor{kwgray}{RGB}{105,105,105}
\definecolor{kwred}{RGB}{150,40,40}
\definecolor{kwgreen}{RGB}{30,95,60}
\definecolor{kwboxbg}{RGB}{247,246,240}

% ---------- 标题样式 ----------
\usepackage{titlesec}
\titleformat{\section}{\Large\bfseries\color{kwblue}}{第\,\thesection\,讲}{0.6em}{}
% switch to plain appendix headings (A, B, C ... instead of 第 N 讲)
\newcommand{\appendixstyle}{%
  \titleformat{\section}{\Large\bfseries\color{kwblue}}{\thesection}{0.8em}{}%
}
\titleformat{\subsection}{\large\bfseries\color{kwblue!85!black}}{\thesubsection}{0.6em}{}
\titleformat{\subsubsection}{\normalsize\bfseries\color{kwblue!70!black}}{\thesubsubsection}{0.5em}{}
\titlespacing*{\section}{0pt}{2.2ex plus .4ex}{1.2ex}
\titlespacing*{\subsection}{0pt}{1.6ex plus .3ex}{0.8ex}

% ---------- 页眉页脚 ----------
\pagestyle{fancy}
\fancyhf{}
\fancyhead[LE,RO]{\small\color{kwgray}\thepage}
\fancyhead[LO]{\small\color{kwgray}Kapustin--Witten 理论导读}
\fancyhead[RE]{\small\color{kwgray}\nouppercase{\leftmark}}
\renewcommand{\headrulewidth}{0.3pt}
\setlength{\headheight}{14pt}
\setlength{\emergencystretch}{2.2em}

% ---------- 定理类 ----------
\theoremstyle{definition}
\newtheorem{definition}{定义}[section]
\newtheorem{example}[definition]{例}
\newtheorem{remark}[definition]{评注}
\newtheorem{exercise}[definition]{习题}
\theoremstyle{plain}
\newtheorem{theorem}[definition]{定理}
\newtheorem{proposition}[definition]{命题}
\newtheorem{lemma}[definition]{引理}
\newtheorem{corollary}[definition]{推论}
\newtheorem{conjecture}[definition]{猜想}
\renewcommand{\proofname}{证明}
\renewcommand{\thedefinition}{\arabic{section}.\arabic{definition}}

% ---------- 讲义专用盒子 ----------
\usepackage[most]{tcolorbox}
\newtcolorbox{goals}{enhanced,breakable,colback=kwblue!4,colframe=kwblue!55,
  boxrule=0.5pt,left=6pt,right=6pt,top=4pt,bottom=4pt,
  fonttitle=\bfseries\small,title={本讲目标}}
\newtcolorbox{follow}{enhanced,breakable,colback=kwgray!6,colframe=kwgray!45,
  boxrule=0.4pt,left=6pt,right=6pt,top=3pt,bottom=3pt,
  fonttitle=\bfseries\small,title={跟读提示（原文页段）}}
\newtcolorbox{keybox}{enhanced,breakable,colback=yellow!6,colframe=yellow!45!orange,
  boxrule=0.5pt,left=6pt,right=6pt,top=4pt,bottom=4pt,
  fonttitle=\bfseries\small,title={本讲核心}}
\newtcolorbox{checkbox}{enhanced,breakable,colback=kwgreen!4,colframe=kwgreen!45,
  boxrule=0.4pt,left=6pt,right=6pt,top=3pt,bottom=3pt,
  fonttitle=\bfseries\small,title={自检问题}}

% ---------- 数学宏 ----------
\newcommand{\MH}{\mathcal{M}_{H}}
\newcommand{\MC}{\mathcal{M}(G,C)}
\newcommand{\Ldual}{{}^{L}G}
\newcommand{\LdualC}{{}^{L}G_{\mathbb{C}}}
\newcommand{\gC}{\mathfrak{g}_{\mathbb{C}}}
\newcommand{\CC}{\mathbb{C}}
\newcommand{\PP}{\mathbb{P}}
\newcommand{\ZZ}{\mathbb{Z}}
\newcommand{\RR}{\mathbb{R}}
\newcommand{\HH}{\mathbb{H}}
\newcommand{\dd}{\mathrm{d}}
\newcommand{\ii}{\mathrm{i}}
\newcommand{\ee}{\mathrm{e}}
\newcommand{\Tr}{\operatorname{Tr}}
\newcommand{\Dmod}{\mathcal{D}}
\newcommand{\kw}{K\!W}
\newcommand{\src}[1]{\textup{\textsf{[#1]}}}
\newcommand{\term}[2]{#1（#2）}

% ---------- 引用键（附录 C 提供 bibitem） ----------
% \cite{kw06} \cite{w08} \cite{fr09} \cite{gw06} \cite{bz04} \cite{set13} \cite{ht03}
